% ---------------------------------------------------------------
\section{Implementation in the BX3 Framework}
\label{sec:bailout-implementation}
% ---------------------------------------------------------------

The BX3 Framework's upstream accountability guarantee states that every
unresolved exception in a node $n$ must reach its designated human anchor
$h(n)$ in finite time, and that the path traversed and the human's directive
must be permanently attributable.  This guarantee is a structural property of
the framework's architecture: it follows from the immutability of the parent
pointer $\alpha(n)$, the finiteness of the spawn tree depth, and the
designation of human anchors at node initialization.  As a structural
property, it holds under the assumption that the framework's enforcement
mechanisms operate correctly.  The Bailout Protocol is the mechanism that
makes this structural guarantee operationally compulsory --- it is the runtime
execution engine that converts the static guarantee into a deterministic,
cryptographically attested, and architecturally enforced propagation event.
The protocol is not an optional extension layered onto the framework; it is
the translation layer between the framework's structural commitments and the
runtime execution trace that constitutes the definitive accountability record
of every exception that occurs in a BX3 node.

The protocol's implementation spans all three BX3 layers, with each layer
fulfilling a distinct and non-overlapping structural role.  The
purposeLayer{} layer defines and anchors the human principal; the \boundsLayer{}
layer houses the Bounds Engine that detects and classifies constraint
violations and emits the trigger condition; the \factLayer{} layer enforces the
state machine transitions, records every event in the Forensic Ledger, and
provides the enforcement substrate that makes compliance architecturally
compulsory rather than operationally discretionary.  This separation of
concerns is not merely organizational --- it is the mechanism by which the
upstream accountability guarantee becomes structurally enforced.

% ---------------------------------------------------------------
\subsection{Purpose Layer: Human Accountability Anchor}
\label{sec:bailout-purpose-layer}

The \purposeLayer{} layer encodes the goal state and terminal success criteria for
the agent's current task and serves as the exclusive terminus for any
unresolved exception in a BX3 node.  The human anchor $h(n)$ is provisioned
at node initialization and stored in \factLayer{}-managed memory.  No machine
actor at any plane P0 through P6 may modify, redirect, or substitute the
human-root identifier.  This immutability is a structural property of the
\purposeLayer{} layer's function, established at provisioning and invariant
throughout the node's operational lifetime.

When the Bailout Protocol state machine transitions to
\texttt{HUMAN\_ACKNOWLEDGED}, the \purposeLayer{} layer receives the escalated
exception with its complete diagnostic context chain assembled by the
escalation algorithm (Algorithm~\ref{alg:escalation}).  The layer receives:
the originating node identifier; the trigger condition's full sub-components
(bound signal, resolution state, confidence value); every resolution attempt
that failed at each intermediate node; the Bounds Engine's reasoning trace at
each hop; and the timeout states that drove the propagation.  This
completeness is enforced by the \factLayer{} layer's pre-commit hook.  No machine
actor may truncate, aggregate, redact, or substitute any element of the
diagnostic context before it reaches the \purposeLayer{} layer.

The \purposeLayer{} layer's role in the protocol is strictly bounded.  It does not
monitor the state machine, does not initiate escalation, and does not
evaluate whether the trigger condition has been satisfied.  It receives
exceptions that have already traversed the full propagation chain and issues
one of three binding directives: \texttt{ACK} (continue with current
parameters, returning the node to \texttt{IDLE}), \texttt{RESOLVE} (execute
a specific resolution prescribed by the human principal, entering
\texttt{RESOLVED}), or \texttt{REJECT} (disallow the proposed action and
enter a quiescent error state, also entering \texttt{RESOLVED}).  The
directive is recorded in the authorization ledger via the \factLayer{} layer before
it takes effect.  The credential binding on the $\sigma_{\text{resolve}}$
event ensures that only $h(n)$ can emit it; no machine actor at any plane
possesses the required credential.  The \purposeLayer{} layer is the structural
realization of the upstream accountability guarantee's terminal commitment:
the exception reaches \emph{that specific human}, no one else, with no
substitutes permitted.

% ---------------------------------------------------------------
\subsection{Bounds Engine: Triggering the Bail-Out}
\label{sec:bailout-bounds-engine}

The \boundsLayer{} layer defines the provisioned operating range $R(n)$ within
which the agent is authorized to act.  The Bounds Engine is the computational
core of this layer: it evaluates whether proposed actions fall within $R(n)$,
executes resolution procedures, produces candidate resolutions, and emits the
trigger condition that activates the Bailout Protocol.

The primary trigger condition TC (Equation~\ref{TC}) fires when three
conditions hold simultaneously: the observable input $x$ falls outside $R(n)$;
$x$ has not been previously resolved in the current operational session; and
the Bounds Engine's self-assessed confidence in its ability to resolve $x$
falls below the provisioned threshold $\tau(n)$.  The Bounds Engine exercises
no discretion over this evaluation --- it performs a deterministic read against
the operating-range definition and the confidence function, both provisioned at
node initialization and immutable at runtime.  The trigger condition either
fires or it does not; the result drives the state machine transition directly,
without an intervening machine-actor gate.

When TC fires, the state machine transitions from \texttt{IDLE} to
\texttt{BOUNDED} (T1, Table~\ref{tab:bp-transitions}), and the Bounds Engine
begins a time-bounded resolution attempt with timeout $T_{\text{bounds}}$.
During \texttt{BOUNDED}, all resolution paths within current authority remain
open; the Bounds Engine may attempt multiple strategies in sequence or in
parallel, each submission reviewed by the \factLayer{} layer against the
authorization ledger.  If the \factLayer{} layer approves a resolution before
$T_{\text{bounds}}$ expires, the state machine returns to \texttt{IDLE} (T2).
If $T_{\text{bounds}}$ expires without approval, the state machine
transitions to \texttt{BAIL\_PENDING} (T3) and the Bounds Engine emits
\texttt{bailout\_signal}.

The transition to \texttt{BAIL\_PENDING} is deterministic and non-negotiable.
No machine actor may suppress the signal, extend $T_{\text{bounds}}$, or
substitute a lower-priority resolution path.  The \factLayer{} layer's pre-commit
hook enforces this: any attempt to write a resolution approval after
$T_{\text{bounds}}$ has expired is rejected, and the protocol violation is
recorded in the Forensic Ledger.  This architectural enforcement is the
structural distinction between the Bailout Protocol and discretionary escalation
interfaces, where the agent decides whether and when to invoke human oversight.

The secondary trigger pathway handles explicit \texttt{bailout\_signal}
emissions from the Bounds Engine for cases where the engine detects a
condition that, while not satisfying TC strictly, represents a threshold
violation warranting human review.  The state machine routes this signal
through the same transition logic (T9, Table~\ref{tab:bp-transitions}), and the
diagnostic context chain is appended identically.  This ensures the protocol
captures not only formally defined TC events but also the broader class of
operator-discretionary escalation events that are contextually critical even
when they do not meet the strict TC criteria.

The Bounds Engine's projection function $\pi(n, x)$ computes the confidence
score $\text{conf}(n, x) \in [0, 1]$ used in the TC evaluation.  This score
is a deterministic read from the engine's state at evaluation time; no machine
actor may bias the score upward to suppress a warranted escalation.  The
threshold $\tau(n)$ governs the conservatism of the node's escalation
behavior: higher values cause earlier escalation, lower values permit deeper
autonomous reasoning before bailout triggers.

% ---------------------------------------------------------------
\subsection{State Machine as Fact Layer Extension}
\label{sec:bailout-state-machine}

The Bailout Protocol state machine
$\mathcal{B} = (Q, \Sigma, \rightarrow, q_0, F)$ is embedded in the
\factLayer{} layer, not in the \boundsLayer{} layer or any machine-executed process.
This embedding is architecturally deliberate: a state machine executing as
a machine process can be intercepted, modified, or circumvented by a
compromised machine actor; a state machine executing as a \factLayer{}
extension is enforced by the layer's deterministic logic, which operates below
the agent's execution level and cannot be intercepted by any machine actor
above it.  The state machine is not a component that the framework
communicates with --- it is a structural extension of the enforcement layer
itself.

The state machine's transition relation $\rightarrow \subseteq Q \times \Sigma \times Q$
is evaluated by the \factLayer{} layer before any state change is committed.  If
no transition is defined for a $(q, \sigma)$ pair, the machine stalls and the
Bailout Monitor records a protocol violation in the Forensic Ledger.  The
state invariants encode the protocol's core guarantees as properties of the
state machine:

\begin{align}
\text{INV-1}&: \forall n \in \text{nodes},\ \text{state}(n) \in Q \setminus \{\texttt{ESCALATING}\} \iff \neg\text{active\_bailout}(n). \tag{BP-INV-1} \\[4pt]
\text{INV-2}&: \text{state}(n) = \texttt{BAIL\_PENDING} \implies \text{active\_bailout}(n) \land \text{transition}(\texttt{BAIL\_PENDING}, \sigma_{\text{timeout}}) = \texttt{ESCALATING}. \tag{BP-INV-2} \\[4pt]
\text{INV-3}&: \text{state}(n) = \texttt{RESOLVED} \implies \exists d \in \{\texttt{ACK}, \texttt{RESOLVE}, \texttt{REJECT}\} \land \text{directive}(n) = d. \tag{BP-INV-3}
\end{align}

INV-2 is particularly significant: it encodes the no-machine-actor-discretion
property as a state machine invariant.  The transition from
\texttt{BAIL\_PENDING} to \texttt{ESCALATING} is forced by the state
machine's deterministic transition relation and enforced by the Bailout Monitor
without machine-actor intervention.  The \factLayer{} layer's pre-commit hook
ensures this transition cannot be intercepted or redirected.

Timeout handling --- $T_{\text{bounds}}$, $T_{\text{prop}}$, $T_{\text{cap}}$,
and $T_{\text{human}}$ (Table~\ref{tab:bp-timeouts}) --- is enforced by the
\factLayer{} layer, not by the \boundsLayer{} layer or any machine actor.  This prevents
a compromised machine actor from artificially extending timeouts to suppress
warranted escalation.  The Pathway Health Registry, maintained by the Bailout
Monitor, tracks the functional status of each parent pointer $\alpha(n)$ for
all nodes.  A pathway is marked \texttt{degraded} when \texttt{SendToParent}
fails $R_{\max}$ times and is excluded from pathway selection until a health
check succeeds.  Degraded pathways are recorded in the Forensic Ledger,
ensuring the failure is attributable rather than silent.

The state machine's embedded position in the \factLayer{} layer ensures that the
Bailout Protocol's guarantees are not contingent on the correct behavior of
the \boundsLayer{} layer or any other machine actor.  The \factLayer{} layer enforces
the protocol's transition relation regardless of what any machine actor
attempted --- it is the structural enforcement substrate that makes the
upstream accountability guarantee architecturally compulsory rather than
operationally optional.

% ---------------------------------------------------------------
\subsection{Fact Layer: Enforcement and Forensic Ledger Recording}
\label{sec:bailout-fact-layer}

The \factLayer{} layer is the enforcement and recording substrate for the Bailout
Protocol.  It maintains the Forensic Ledger --- an immutable, SHA-256
hash-chained record of all state transitions, trigger conditions, propagation
events, and human directives across the protocol's lifetime.  The Forensic
Ledger serves two distinct but related functions: it is the authoritative
runtime state store for the Bailout Protocol's state machine, and it is the
enforcement substrate for the framework's accountability properties.  The
pre-commit hooks for the Bailout Protocol, Sandbox Gate, and Spatial Firewall
all query the Forensic Ledger to determine the current state of their
respective state machines before any state change is committed.

The Bailout Protocol generates the following entry types in the Forensic
Ledger, each written atomically by the \factLayer{} layer before the corresponding
state transition is committed:

\begin{itemize}\itemsep=0pt
\item[\texttt{bailout-trigger}] Originating node, exception identifier,
  trigger-condition sub-components (bound signal, resolution state,
  confidence value), and priority function output $\pi(n, x)$.
\item[\texttt{bounds-resolution-attempt}] Each resolution procedure submitted
  by the Bounds Engine during \texttt{BOUNDED} state, including the
  authorization ledger entry reference and the \factLayer{} approval or
  rejection outcome.
\item[\texttt{bailout-signal}] Explicit \texttt{bailout\_signal} emissions,
  including the Bounds Engine's stated reason.
\item[\texttt{timeout-bounds-expired}] Expiration of $T_{\text{bounds}}$ and
  absence of a \factLayer{}-approved resolution at expiration.
\item[\texttt{escalation-sent}] Each invocation of \texttt{SendToParent},
  including the pathway selected, the payload hash, the retry count, and the
  acknowledgment status.
\item[\texttt{pathway-exhausted}] Exhaustion of all retry pathways and the
  terminal state reached.
\item[\texttt{human-ack}] Human anchor's acknowledgment of the escalated
  exception, acknowledgment timestamp, and pathway used.
\item[\texttt{human-directive}] Binding directive issued by $h(n)$ --- one of
  \texttt{ACK}, \texttt{RESOLVE}, or \texttt{REJECT} --- and the
  corresponding authorization ledger entry hash.
\end{itemize}

The Forensic Ledger entry for a Bailout Protocol escalation constitutes the
definitive accountability record for that exception.  It captures the complete
provenance of the exception from the triggering input through every
intermediate node's resolution attempt to the human principal's directive.
No machine actor may modify, delete, or suppress any entry; the \factLayer{} layer
enforces the record rather than interprets it.  This is the architectural
distinction between the \factLayer{} layer and a conventional audit log: the
\factLayer{} layer does not merely store events --- it actively prevents invalid
state transitions, making it an enforcement mechanism rather than a passive
recorder.

The SHA-256 hash chaining ensures tamper-evidence: each entry includes the
hash of the preceding entry, making any post-hoc modification detectable by
recalculating the chain.  An external auditor --- regulator, customer,
investor, or court --- can reconstruct the complete exception history of any
node from the cryptographically chained Forensic Ledger entries and verify
that every unresolved exception reached $h(n)$, that every human directive
was recorded, and that no machine actor suppressed, redirected, or substituted
for the human principal.  The guarantee is not merely asserted --- it is
evidenced.

Completeness of the Forensic Ledger is enforced by Invariant I4: every state
transition in the Bailout Protocol's state machine produces at least one
Forensic Ledger entry before the transition is committed.  This is not an
operational convention; it is enforced by the \factLayer{} layer's pre-commit
hook, which rejects any state transition that has not produced its required
entry.  The combination of hash chaining and completeness enforcement means
that the Forensic Ledger is both tamper-evident and exhaustive: any gap in
the chain is itself a detectable and attributable protocol violation.

% ---------------------------------------------------------------
\subsection{Runtime Expression of the Upstream Accountability Guarantee}
\label{sec:bailout-uag}

The upstream accountability guarantee exists as a structural property of the
BX3 Framework until an exception triggers the Bailout Protocol, at which point
the protocol converts it into an immutable, evidenced execution trace.  This
translation from static guarantee to runtime event is the protocol's defining
function: it is the mechanism that makes the upstream accountability guarantee
operationally real, converting the framework's architectural commitments into
concrete, attributable outcomes.

Without the Bailout Protocol, the upstream accountability guarantee would be a
structural aspiration with no enforcement mechanism.  The human anchor $h(n)$
would be the designated terminus in principle, but nothing would compel
exceptions to reach it.  Machine actors could suppress, redirect, or silently
ignore exception states, and the guarantee would hold only for systems whose
machine actors chose to comply.  The Bailout Protocol converts the guarantee
from a structural aspiration into an architecturally enforced compulsion: the
state machine's deterministic transition relation, the \factLayer{} layer's
pre-commit hook, and the immutable parent-pointer architecture together ensure
that no unresolved exception state can exist in a BX3 node without propagating
to $h(n)$.

The Forensic Ledger entries generated by the protocol provide the retrospective
proof of this guarantee.  Each entry is written before the corresponding state
transition commits, forming a complete and attributable chain from the
triggering exception to the human principal's directive.  An external auditor
can verify the guarantee's fulfillment by examining the chain: every
\texttt{bailout-trigger} entry must have a corresponding \texttt{human-directive}
entry, with all intermediate propagation steps recorded.  The chain is
verifiable in both directions --- forward from trigger to resolution, and
backward from any recorded directive to its originating trigger.

The Bailout Protocol therefore occupies a unique position in the BX3
Framework: it is the mechanism that makes the upstream accountability guarantee
operationally real.  It is the translation layer between the framework's
structural commitments (three-layer separation, immutable parent pointers,
human anchor designation) and the runtime execution trace (state transitions,
timeout events, propagation paths, human directives) that constitutes the
definitive accountability record of the framework.  The protocol is not an
extension or an optional feature --- it is the runtime expression of the
upstream accountability guarantee in exactly the same sense that the state
machine is the runtime expression of the transition relation.